\documentclass[journal]{IEEEtran}

\usepackage{amsmath,amssymb}
\usepackage{booktabs}
\usepackage{array}
\usepackage{cite}
\usepackage{graphicx}
\usepackage{multirow}
\usepackage{threeparttable}
\usepackage{url}

% Keep full-width tables visually attached to the surrounding text.
\setlength{\textfloatsep}{8pt plus 2pt minus 2pt}
\setlength{\dbltextfloatsep}{8pt plus 2pt minus 2pt}
\setlength{\floatsep}{6pt plus 2pt minus 2pt}
\setlength{\dblfloatsep}{8pt plus 1pt minus 1pt}

\newcommand{\auroc}{\mathrm{AUROC}}
\newcommand{\auprc}{\mathrm{AUPRC}}
\newcommand{\ece}{\mathrm{ECE}}
\newcolumntype{L}[1]{>{\raggedright\arraybackslash}p{#1}}

\title{Supplementary Material for ``Multimodal COVID-19 Respiratory Sound Models: Temporal Stability and External Cough Transfer''}
\author{Nishant Harkut and Ritu Tiwari}

\begin{document}
\bstctlcite{BSTcontrol}
\maketitle

\section{Purpose and Analysis Scope}

This supplement records implementation details and secondary analyses needed for reproducibility but omitted from the main article for readability. No result in this document was used to alter a source-to-target model after inspecting COUGHVID performance. The primary analysis is the model- and modality-matched Coswara-cough to COUGHVID-cough transfer. Learned-representation transfer and calibration are supporting analyses. Temporal, metadata, feature-stability, overlap, subgroup, and incremental-value analyses are secondary or exploratory. The study was not preregistered.

\section{Dataset Endpoints and Inclusion Rules}

\subsection{Coswara}

Coswara contains multiple voluntary respiratory recordings and participant-supplied health fields \cite{bhattacharya2023coswara}. Positive asymptomatic, mild, and moderate categories were mapped to the positive class. Healthy participants were mapped to the negative class. Recovered, exposed, other respiratory illness, under-validation, and unresolved records were excluded from supervised endpoints. The resulting known-label source cohort contained 2,114 participants, including 681 positives and 1,433 negatives.

\subsection{COUGHVID}

COUGHVID is a cough-only crowdsourced resource with heterogeneous participant-supplied and annotated fields \cite{orlandic2021coughvid}. The predetermined external table retained 8,331 unique recording UUIDs and used the semi-supervised \texttt{status\_SSL} field produced by combining user labels with models trained on sparse expert annotations \cite{orlandic2023semisupervised}. Values designated COVID-19 or positive were mapped to positive, and values designated healthy or negative to negative. The table contained 285 positive-labeled and 8,046 negative-labeled recordings, for positive-label prevalence 0.0342. The public table provides one row per recording rather than a persistent person identifier, so repeated submissions by one individual cannot be grouped. The two datasets do not provide identical clinical reference standards. External performance therefore measures recording-level transport and agreement between operational dataset endpoints, not accuracy against two harmonized polymerase-chain-reaction reference cohorts.

\begin{table*}[t]
\caption{Label and Evaluation Boundaries}
\label{tab:s_label_boundaries}
\centering
\footnotesize
\begin{tabular}{L{2.1cm}L{3.2cm}L{3.2cm}L{5.6cm}}
\toprule
Component & Included positive endpoint & Included negative endpoint & Boundary on interpretation \\
\midrule
Coswara & Available positive asymptomatic, mild, and moderate categories & Healthy category & Self-reported operational categories. Unresolved, recovered, exposed, and other respiratory-illness categories excluded \\
COUGHVID & Available target positive status after predetermined mapping and deduplication & Available target negative status & Cough-only collection with different recruitment and label provenance \\
External transfer & Coswara-trained cough score evaluated against the COUGHVID endpoint & Same & Tests dataset transport under combined endpoint and acquisition change. Not a harmonized clinical diagnostic study \\
Multimodal fusion & Coswara participant endpoint & Coswara participant endpoint & Source-only because COUGHVID has no matched breath or speech modalities \\
\bottomrule
\end{tabular}
\end{table*}

\section{Audio Processing and Quality Control}

Files were decoded to mono floating-point waveforms and resampled to 16~kHz. The audit identified 23,539 usable recordings, 597 mostly silent recordings, 390 corrupt recordings, and 186 short recordings. Quality flags were determined before model fitting and describe file usability rather than disease status.

For the engineered representation, silence trimming used a 35-dB threshold. An energy-based event detector used 25-ms frames and 10-ms hops. Frames above 25\% of the maximum frame energy were retained with a 120-ms margin. The selected event was peak-normalized and center-cropped or zero-padded to 4~s for cough and 5~s for breath and speech. Event start, event end, active duration, and active-audio fraction were retained as acquisition-context descriptors. Their presence is reported explicitly because timing can reflect recording behavior as well as physiology.

\begin{table*}[t]
\caption{Deterministic Audio Processing Settings}
\label{tab:s_audio_settings}
\centering
\footnotesize
\begin{tabular}{L{3.0cm}L{3.0cm}L{7.3cm}}
\toprule
Stage & Setting & Purpose and boundary \\
\midrule
Decoding & Mono, floating point & Standardized channel layout with unreadable files flagged \\
Sampling rate & 16~kHz & Common rate for engineered and learned branches \\
Silence trim & 35~dB & Removes low-energy leading and trailing regions \\
Energy detector & 25-ms frame, 10-ms hop, 25\% peak-energy threshold & Locates the active event using an auditable deterministic rule \\
Event margin & 120~ms & Retains local context around detected activity \\
Amplitude & Peak normalization & Reduces gross recording-level scale differences \\
Fixed duration & 4~s for cough and 5~s for breath or speech & Center crop or zero pad to a stable analysis window \\
Timing descriptors & Event boundaries, active duration, active fraction & Included and disclosed as context-sensitive descriptors \\
\bottomrule
\end{tabular}
\end{table*}

\section{Acoustic Representations}

\subsection{Engineered Representation}

The merged table contained three families rather than only openSMILE features. The first was a broad acoustic bank computed from the processed waveform using librosa-based extraction \cite{mcfee2015librosa}. The other two were openSMILE ComParE 2016 and IS10 paralinguistic sets \cite{eyben2010opensmile,schuller2010is10,schuller2016compare}. After identifier and context columns, the merged table contained 10,147 columns before global selection.

\begin{table*}[t]
\caption{Engineered Feature Families}
\label{tab:s_features}
\centering
\footnotesize
\begin{tabular}{L{3.0cm}L{4.2cm}L{6.2cm}}
\toprule
Family & Examples & Information summarized \\
\midrule
Cepstral and mel & 40 MFCC trajectories, first and second differences, 64 log-mel bands & Spectral envelope, time-varying changes, and band-energy distribution \\
Spectral shape & Centroid, bandwidth, roll-off, flatness, contrast & Frequency concentration, spread, edge, tonality, and local contrast \\
Energy and voicing proxies & RMS energy, zero-crossing rate & Amplitude dynamics and rapid sign-change behavior \\
Pitch-related and tonal & Chroma and tonal centroid features & Pitch-class and harmonic organization \\
Temporal rhythm & Tempogram and trajectory summaries & Modulation and rate of acoustic change \\
Acquisition context & Detected event start, end, duration, active fraction & Recording and event timing, not interpreted as disease-specific physiology \\
openSMILE ComParE 2016 & Challenge functionals over low-level descriptors & Broad paralinguistic reference representation \\
openSMILE IS10 & Interspeech 2010 paralinguistic functionals & Compact reference representation \\
\bottomrule
\end{tabular}
\end{table*}

Constant columns were removed. LightGBM gain importances were computed separately on the training rows for each available source modality. Each modality-specific importance vector was divided by its maximum, and the normalized vectors were averaged. The 800 highest-ranked columns were frozen before external scoring. Within each conventional classifier pipeline, a training-fitted variance filter was followed by an ANOVA percentile filter. Boosted-tree pipelines retained 80\%. RBF-SVC retained 60\% after standardization. No COUGHVID label or performance measure influenced selection.

\subsection{Learned Representations}

WavLM-base-plus provided a self-supervised transformer check \cite{wang2022wavlm}. Up to four 3-s cough segments with 50\% overlap were processed per recording. Segment probabilities were averaged per Coswara participant and per COUGHVID recording UUID. The source test and external target contained 316 participants and 8,331 recording UUIDs. The CNN--BiGRU branch used one 64-band log-mel spectrogram per recording, convolutional feature extraction, and a bidirectional gated recurrent unit for temporal aggregation. Its source test and external target contained 636 recordings and 8,331 recording UUIDs. These branches used different training recipes and analysis units, so they are architecture checks rather than members of the controlled four-model comparison.

\begin{table*}[t]
\caption{Settings for Supporting Learned-Representation Checks}
\label{tab:s_learned_settings}
\centering
\footnotesize
\begin{tabular}{L{2.4cm}L{5.2cm}L{5.8cm}}
\toprule
Branch & Representation and classifier & Optimization and selection \\
\midrule
WavLM-base-plus & Mean-pooled encoder states followed by layer normalization, 256-unit hidden layer, GELU activation, and one output logit. The encoder feature extractor was frozen and the top four transformer layers were trainable & At most 8 epochs, physical batch 8, gradient accumulation 4, encoder and head learning rates $2\times10^{-5}$ and $10^{-4}$, AdamW weight decay 0.01, class-weighted binary cross-entropy, cosine schedule with 10\% warm-up, and patience 3 on participant-level validation AUROC \\
CNN--BiGRU & Three $3\times3$ convolutional layers with 32, 64, and 64 channels, frequency-axis pooling after the first two layers, one bidirectional GRU with 96 hidden units per direction, mean temporal pooling, 128-unit hidden layer, and dropout 0.35 & At most 8 epochs, batch 16, AdamW learning rate $10^{-3}$ and weight decay $10^{-4}$, class-weighted binary cross-entropy, training-only noise and time-frequency masking, and patience 3 on recording-level validation AUROC \\
\bottomrule
\end{tabular}
\end{table*}

\section{Conventional Models and Imbalance Handling}

All transformations, feature filters, scaling, oversampling, and model fitting were restricted to the corresponding training partition. Synthetic minority oversampling used three nearest neighbors for the boosted-tree pipelines \cite{chawla2002smote}. Hyperparameters were frozen before external scoring.

\begin{table*}[t]
\caption{Fixed Conventional Model Settings}
\label{tab:s_models}
\centering
\footnotesize
\begin{tabular}{L{2.2cm}L{2.0cm}L{9.3cm}}
\toprule
Model & Feature percentile & Main settings \\
\midrule
LightGBM & 80\% & 900 estimators, learning rate 0.025, 31 leaves, minimum child size 20, row and feature subsampling 0.85, $L_2=2$, and training-only SMOTE \cite{ke2017lightgbm} \\
XGBoost & 80\% & 800 estimators, depth 3, learning rate 0.025, row and feature subsampling 0.85, minimum child weight 2, $L_2=2$, and training-only SMOTE \cite{chen2016xgboost} \\
CatBoost & 80\% & 900 iterations, depth 5, learning rate 0.025, and training-only SMOTE \cite{prokhorenkova2018catboost} \\
RBF-SVC & 60\% & Standard scaling, $C=2$, class balancing, and probabilities from three-fold sigmoid calibration \cite{platt1999probabilistic,niculescu2005calibration} \\
\bottomrule
\end{tabular}
\end{table*}

\section{Participant Aggregation and Multimodal Fusion}

For participant $i$, modality $m$, and available recordings $r=1,\ldots,R_{im}$, recording probabilities were averaged:
\begin{equation}
p_{im}=\frac{1}{R_{im}}\sum_{r=1}^{R_{im}}p_{irm}.
\end{equation}
Only after participant aggregation were modalities fused. Uniform fusion used their arithmetic mean. Validation-weighted fusion assigned
\begin{equation}
\widetilde w_m=\max\!\left(\auprc_m^{(val)}-0.5,0.01\right),\qquad
w_m=\frac{\widetilde w_m}{\sum_j\widetilde w_j},
\end{equation}
and combined $p_i=\sum_mw_mp_{im}$. Stacked fusion fitted class-balanced logistic regression to validation-set modality probabilities with a maximum of 2,000 iterations. Thus, no final modality weight was hard-coded. Binary thresholds maximized validation balanced accuracy and were then frozen for source test evaluation.

The selected participant-split source model used cough and speech in a validation-trained logistic stack. It reached AUROC 0.897 and AUPRC 0.863 on 314 complete-case source-test participants. This is a source multimodal result. The cough-only COUGHVID experiment does not test this fusion.

\section{Evaluation Protocols}

\begin{table*}[t]
\caption{Protocol Definitions and What Each Can Establish}
\label{tab:s_protocols}
\centering
\footnotesize
\begin{tabular}{L{2.6cm}L{3.7cm}L{3.4cm}L{3.8cm}}
\toprule
Protocol & Construction & Intended question & Interpretation boundary \\
\midrule
Participant split & Stratified 70/15/15 source split with all recordings from each participant kept together & Can the selected workflow discriminate held-out participants from the same source collection? & Does not test a new collection process \\
Time-stratified participant split & Participant-disjoint source sets distributed across collection time with similar prevalence & Is source performance sensitive to a different time-aware assignment? & Not future deployment. Complete-case denominator depends on modalities \\
Early-to-late stress test & Train through 2020-12-04, validate through 2021-07-29, and test later recordings & Is the workflow stable under a large calendar and prevalence change? & Retrospective and reuses the broader source-selected representation \\
External transfer & Train and tune on Coswara, then score untouched COUGHVID & Does cough discrimination transport to an independent dataset endpoint? & Source and target label mechanisms are not harmonized \\
Target recalibration & Fit on 2,082 target recording UUIDs and evaluate on 6,249 disjoint UUIDs & Is failure mainly probability-scale error or ranking error? & Supervised target adaptation, not an unsupervised deployment procedure \\
\bottomrule
\end{tabular}
\end{table*}

These protocols answer different questions and are not a monotonic validation ladder. In particular, the internal multimodal endpoint cannot be directly subtracted from the external cough-only endpoint.

\section{Metrics and Statistical Procedures}

Discrimination was measured using AUROC and AUPRC. Threshold-dependent metrics were balanced accuracy, sensitivity, specificity, F1 score, and precision. Probability quality was summarized using negative log likelihood, Brier score \cite{brier1950verification}, and ECE \cite{vancalster2019calibration}. For observations $y_i\in\{0,1\}$ and probabilities $p_i$,
\begin{equation}
\mathrm{Brier}=\frac{1}{n}\sum_{i=1}^{n}(p_i-y_i)^2.
\end{equation}
For bins $b=1,\ldots,B$,
\begin{equation}
\ece=\sum_{b=1}^{B}\frac{n_b}{n}\left|\overline p_b-\overline y_b\right|.
\end{equation}

Coswara participants and COUGHVID recording UUIDs are independent across datasets. The interval for $\Delta=\auroc_{source}-\auroc_{target}$ was obtained from 2,000 independent bootstrap replicates. Each replicate resampled source participants and target recording UUIDs separately, recomputed both AUROCs, and stored their difference. A base random seed of 42 controlled reproducibility, and percentile intervals used the 2.5th and 97.5th quantiles. This is not subtraction of two marginal intervals. DeLong tests were used only when two scores were available for the same participants \cite{delong1988roc}. AUPRC and precision were interpreted relative to endpoint prevalence \cite{saito2015auprc}.

Operating-point, subgroup, context-control, and related extensions were exploratory. Their intervals and $p$-values were not multiplicity-adjusted to support separate primary claims. Because the COUGHVID endpoint is a semi-supervised pseudo-label, external sensitivity, specificity, and precision characterize label agreement rather than clinical performance.

\section{Detailed External Transfer Results}

\begin{table*}[t]
\caption{Controlled Cough-Only Source-to-Target Transfer}
\label{tab:s_external}
\centering
\footnotesize
\begin{tabular}{lrrrrrr}
\toprule
Model & Source AUROC & Target AUROC & $\Delta$ AUROC & 95\% CI for $\Delta$ & Source AUPRC & Target AUPRC \\
\midrule
CatBoost & 0.849 & 0.531 & 0.318 & [0.256, 0.376] & 0.788 & 0.040 \\
LightGBM & 0.853 & 0.543 & 0.310 & [0.247, 0.366] & 0.797 & 0.040 \\
RBF-SVC & 0.868 & 0.523 & 0.345 & [0.291, 0.398] & 0.812 & 0.037 \\
XGBoost & 0.850 & 0.532 & 0.318 & [0.262, 0.371] & 0.780 & 0.039 \\
\bottomrule
\end{tabular}
\end{table*}

The lower confidence bound for every source-to-target decline exceeded 0.247. Because modality, selected representation, model family, and transfer direction were fixed, this is the most tightly controlled comparison in the study.

\begin{table*}[t]
\vspace{4pt}
\caption{Learned-Representation External Checks}
\label{tab:s_deep}
\centering
\footnotesize
\begin{tabular}{lrrrrrr}
\toprule
Representation & Source $n$ & Target $n$ & Source AUROC & Target AUROC & Source AUPRC & Target AUPRC \\
\midrule
WavLM-base-plus pooled cough & 316 & 8,331 & 0.812 & 0.484 & 0.734 & 0.032 \\
CNN--BiGRU spectrogram model & 636 & 8,331 & 0.737 & 0.548 & 0.574 & 0.044 \\
\bottomrule
\end{tabular}
\end{table*}

These results show that the observed transport problem was not confined to one engineered representation. They do not constitute an architecture-controlled comparison because the learned branches used different training procedures and source denominators.

\section{Operational-Label Agreement and Recalibration}

At sensitivity at least 0.90 against \texttt{status\_SSL}, specificity ranged from 0.112 to 0.161 and precision from 0.0347 to 0.0367. Positive-label prevalence was 0.0342, so precision was only marginally above the pseudo-label base rate. This operating point describes agreement with the operational target label. It is not a clinical screening estimate or a deployable target-threshold procedure.

\begin{table*}[t]
\caption{Held-Out COUGHVID Recalibration Example for LightGBM}
\label{tab:s_recalibration}
\centering
\footnotesize
\begin{tabular}{lrrrrr}
\toprule
Mapping & Calibration $n$ & Evaluation $n$ & Evaluation AUROC & Brier score & ECE \\
\midrule
Original probability & 2,082 & 6,249 & 0.539 & 0.0756 & 0.0992 \\
Platt mapping & 2,082 & 6,249 & 0.539 & 0.0330 & 0.00017 \\
Isotonic mapping & 2,082 & 6,249 & 0.532 & 0.0331 & 0.00027 \\
\bottomrule
\end{tabular}
\end{table*}

Calibration maps were fitted on the 2,082-recording calibration subset and evaluated on a disjoint 6,249-recording subset. Platt scaling reduced measured calibration error while preserving AUROC at 0.539. Isotonic regression also reduced calibration error but did not improve AUROC. Poor agreement with the operational endpoint is therefore not explained by threshold or probability-scale mismatch alone.

\section{Temporal and Repeated-Seed Analyses}

The selected participant-split fusion remained stable across four seeds, with mean AUROC $0.895\pm0.003$. The retrospective early-to-late endpoint had mean AUROC $0.691\pm0.007$ across three seeds. Train, validation, and test prevalence was 0.079, 0.569, and 0.804. These changes make AUPRC and calibration strongly prevalence-dependent and prevent a simple causal attribution to biological temporal drift.

The reverse late-to-early experiment yielded AUROC 0.920 but AUPRC 0.029, F1 score 0.011, and ECE 0.471. This illustrates why rank discrimination alone is insufficient under extreme base-rate change. It is secondary and is not evidence of useful reverse deployment.

An exploratory LightGBM ranking comparison used an independent cutoff of 13 August 2020. The early and late top-800 lists shared 110 features, giving Jaccard similarity $110/1490=0.074$. The early rows contained 52 positive recordings and the late rows 3,983. The low overlap is consistent with non-stationary associations, but it is entangled with major prevalence and recruitment changes.

\section{Collection-Context and Metadata Controls}

Source labels were predictable from non-audio fields. The full context model reached AUROC 0.964, the symptoms-only model 0.932, and the demographic/protocol model 0.914. After full retraining with shuffled labels, mean AUROCs were approximately 0.503 for all three. The collapse after label permutation is an implementation-leakage sanity check. It does not remove confounding from observed-label data.

\begin{table*}[t]
\caption{Grouped Permutation Attribution for Context Models}
\label{tab:s_context}
\centering
\footnotesize
\begin{tabular}{llrrl}
\toprule
Model & Feature group & Features & Importance share & Highest-importance field \\
\midrule
Full context & Recording protocol & 8 & 0.717 & Recording year \\
Full context & Symptoms or label proxies & 10 & 0.185 & Reported cough symptom \\
Full context & Comorbidity proxies & 8 & 0.052 & Other pre-existing condition \\
Full context & Demographics & 37 & 0.046 & Country indicator \\
Demographic/protocol & Recording protocol & 8 & 0.934 & Recording year \\
Demographic/protocol & Demographics & 37 & 0.066 & Country indicator \\
\bottomrule
\end{tabular}
\end{table*}

These values document label--context association. They do not demonstrate that a particular audio model encoded every metadata field. Direct attribution would require matched or interventional analysis linking audio scores to specific context variables.

\section{Source--Target Separation and Support Diagnostic}

A standardized, class-balanced logistic classifier trained to distinguish source from target reached five-fold out-of-fold AUROC 0.750. It used the 500 common features with highest variance in source-training rows. Mean predicted target-domain probability was 0.399 for source observations and 0.629 for target observations. Using the 5th--95th percentile probability band of source observations as a descriptive reference, 25.2\% of external observations fell outside the source band. This indicates measurable source--target separation, but the one-dimensional probability-band statistic is not a formal proof of causal positivity or full geometric non-overlap.

\section{Incremental Value of Audio}

Participant-aligned analyses compared metadata-only, audio-only, and combined logistic models on complete cases. The symptoms-only plus best aligned CatBoost audio comparison used 61 participants. Metadata AUROC was 0.888, audio AUROC 0.818, and combined AUROC 0.951. The combined-minus-metadata difference was 0.063, with bootstrap interval $[-0.005,0.149]$ and paired DeLong $p=0.104$. For full context plus the selected fusion on 61 participants, metadata AUROC was 0.924 and combined AUROC 0.933. The difference was 0.009, interval $[-0.026,0.040]$, and $p=0.538$. Full context plus a LightGBM feature-fusion score gave a difference of 0.006, interval $[-0.038,0.048]$, and $p=0.735$.

The complete-case denominator is too small to exclude modest incremental value. These analyses show no statistically resolved gain in the available aligned subset. They do not prove equivalence between metadata-only and combined prediction.

\section{Multiplicity and Reporting Discipline}

The four-model cough-only source-to-target comparison is the primary evidence. Deep transfer, calibration, and high-sensitivity operating points support its interpretation. Remaining analyses diagnose stability, context, or plausibility after the core result was established. Exploratory intervals and $p$-values are effect-size descriptors and were not adjusted to claim additional primary findings. This distinction is necessary because the evidence package contains many models, subgroups, thresholds, and post-hoc checks.

\section{Reproducibility Inventory}

The reproducibility package should contain participant-level source predictions where aggregation was performed, recording-level CNN and external predictions, metric tables, bootstrap estimates, calibration-bin tables, protocol and leakage audits, feature-selection summaries, model-selection summaries, deterministic figure-generation code, and a machine-readable experiment manifest. Raw audio remains subject to the respective dataset terms. Public release should not redistribute participant-level fields where source licenses prohibit it.

\bibliographystyle{IEEEtran}
\IEEEtriggeratref{15}
\bibliography{references}

\end{document}
